\project{07}{Consumer 9-DoF plus climate (IKS4A1)}{Raspberry Pi 4 or 3B+ + Nucleo-H7A3ZI-Q + X-NUCLEO-IKS4A1}{consumer MEMS, humidity and the Qvar electrode}

\begin{unique}
Owns consumer MEMS plus SHT40 humidity and the Qvar electrode. Wearable physics, not
industrial high-g.
\end{unique}

\subsection*{Intent}

The IKS4A1 on the Nucleo streams JSON of the LSM6DSV16X and LSM6DSO16IS inertial parts,
the LIS2MDL magnetometer, the LIS2DUXS12 accelerometer, the LPS22DF barometer, the SHT40
humidity sensor and the STTS22H temperature sensor. Linux records. That is the whole
product: a line of JSON per sample period, written to a file that is honest about what was
measured and when.

The teaching point of this lab is a boundary, not a sensor. Every one of those chips sits
on the STM32 I2C bus, not on the Pi I2C bus, so \texttt{i2cdetect} on the Pi finds nothing
at all. Newcomers read that empty grid as a fault and start pulling jumpers. It is not a
fault, it is the architecture: the microcontroller owns the sensor bus, and Linux owns
storage, time and the file format. The only thing that crosses the USB cable is text the
firmware chose to print. P08 uses the same physical stack with the industrial shield, and
P18 reuses this shield when the M7 runs a filter, but the boundary is identical in all
three.

\begin{kit}
Pi 4 or 3B+, Nucleo-H7A3ZI-Q, X-NUCLEO-IKS4A1, USB cable.
\end{kit}

\subsection*{System architecture}

\diagram{p07_arch}{The boundary that defines this lab. Seven sensor parts sit on the STM32
I2C bus behind the microcontroller. The Pi I2C controller is idle and unwired, which is
why \texttt{i2cdetect -y 1} on the Pi returns an empty grid.}

Read the figure from the left. The shield chips answer at 0x6A for the LSM6 inertial part,
0x1E for the LIS2MDL, 0x44 for the SHT40, 0x5C for the LPS22DF and 0x38 for the STTS22H.
Those addresses live on the microcontroller bus. The scanner that proves them is the one
in the firmware, not the one on the Pi. The Pi I2C controller is drawn deliberately as an
empty box: nothing is wired to pins 3 and 5, and nothing should be.

On the Linux side the work is small and should stay small. A CDC ACM character device
delivers lines, and \texttt{tee} writes them to a file while you watch them go past. Any
processing beyond that, including the tilt computation, is a consumer of the file or of
the same stream, never a replacement for it. Keep the raw capture; a derived quantity you
cannot recompute from the record is a claim rather than a measurement.

\subsection*{Wiring and schematic}

\diagram{p07_schematic}{The IKS4A1 on the Arduino UNO headers, with the UM3239 Mode 1
jumper positions that put every sensor on the microcontroller I2C bus. The Pi side of the
drawing carries one USB connector and nothing else.}

\begin{table}[H]\centering\small
\begin{tabular}{p{40mm}p{28mm}p{22mm}p{58mm}}
\toprule
Signal or part & Connector & Position & Note \\
\midrule
X-NUCLEO-IKS4A1 & Nucleo Arduino UNO headers & stacked & One shield only, IKS5A1 stays bagged \\
UM3239 Mode 1, J4 & IKS4A1 jumper J4 & 1-2 & All sensors on the uC I2C bus \\
UM3239 Mode 1, J4 & IKS4A1 jumper J4 & 11-12 & Second position of the same mode \\
UM3239 Mode 1, J5 & IKS4A1 jumper J5 & 1-2 & All sensors on the uC I2C bus \\
UM3239 Mode 1, J5 & IKS4A1 jumper J5 & 11-12 & Second position of the same mode \\
LSM6DSV16X or LSM6DSO16IS & STM32 I2C bus & 0x6A & Inertial, behind the STM32 \\
LIS2MDL & STM32 I2C bus & 0x1E & Magnetometer \\
SHT40AD1B & STM32 I2C bus & 0x44 & Relative humidity \\
LPS22DF & STM32 I2C bus & 0x5C & Barometer \\
STTS22H & STM32 I2C bus & 0x38 & Temperature \\
LIS2DUXS12 & STM32 I2C bus & see the shield manual & The source does not fix this address, confirm in UM3239 and on the silkscreen \\
Qvar electrode & IKS4A1 jumpers & UM3239 Qvar mode & Left off until the IMU stream is clean \\
Nucleo USB & Pi USB-A & any & CDC ACM, the only link to Linux \\
Pi 40-pin header & Pi & empty & No HAT, no flying leads, nothing on pins 3 and 5 \\
\bottomrule
\end{tabular}
\caption{Wiring. Mode 1 is two jumper positions on each of J4 and J5. Nothing in this lab
connects to the Pi header.}
\end{table}

\begin{asciiart}
IKS4A1 on Nucleo Arduino headers.
UM3239 Mode 1 (all sensors on uC I2C):
  J4: 1-2 and 11-12
  J5: 1-2 and 11-12
Nucleo USB --> Pi.
sudo i2cdetect on the *Pi* will NOT see these chips.
They live behind the STM32. That is the point.
\end{asciiart}

Mode 1 is the Linux-friendly setup because it puts every part on one bus that the firmware
owns and can scan in one pass. The other modes on this shield route parts through the
sensor hub or the ISPU inside the LSM6, which hides devices from that scan. Those modes are
worth a session after the JSON is honest, not before.

\subsection*{Bench layout}

\diagram{p07_bench}{Bench layout. The shield is stacked on the Nucleo, one USB cable
reaches the Pi, and the Pi header stays empty. The two stimuli that prove the climate
sensors are a breath toward the SHT40 and a palm held over the board.}

Put the Nucleo where you can reach the shield with a hand and a breath without moving the
cable. Keep the stack away from a laptop fan and away from a mug: both move relative
humidity and temperature on their own and make the acceptance test ambiguous. The Pi can
sit anywhere the USB cable reaches, because it wears no HAT in this lab and has nothing
plugged into its header.

\subsection*{Software design (UML)}

\diagram{p07_uml}{Component view. The seven sensor components sit behind the STM32 I2C
master. The firmware exposes one interface to the outside world, a CDC ACM port carrying
JSON lines, and the Linux recorder consumes exactly that. There is no interface from the
Pi to the sensor bus, which is why the Pi scanner is empty.}

The component diagram is the argument of this chapter drawn once. Every sensor is a
component with a required I2C interface, and the only component that provides that
interface is the STM32 I2C master. The firmware provides a single external interface: JSON
lines over CDC. The Linux recorder requires that interface and provides a file. Nothing in
the picture connects a Linux component to a sensor component, and no wiring change in this
lab can create such a connection.

\subsection*{Data flow (ASCII)}

\begin{asciiart}
  X-NUCLEO-IKS4A1                       Nucleo-H7A3ZI-Q          Pi 4 or 3B+
  +-----------------------+             +-----------------+      +-----------------+
  | LSM6DSV16X / DSO16IS  |--0x6A--\    |                 |      |                 |
  | LIS2MDL               |--0x1E---\   | STM32 I2C       |      | /dev/ttyACM0    |
  | SHT40AD1B             |--0x44----+->| master + driver |      |   |             |
  | LPS22DF               |--0x5C---/   |     |           | USB  |   v             |
  | STTS22H               |--0x38--/    |     v           |=====>| cat | tee       |
  | LIS2DUXS12            |--see UM3239 | one JSON object |  CDC |   p07.jsonl     |
  | Qvar electrode (off)  |             | per 20 ms       |      |   |             |
  +-----------------------+             +-----------------+      |   v             |
                                                                 | tilt from acc   |
                                                                 | and gyro only   |
                                                                 +-----------------+

  Pi i2c-1:  nothing wired, i2cdetect -y 1 is empty.  That is the architecture,
             not a fault.  The scanner that proves the addresses is in the firmware.
\end{asciiart}

\subsection*{Steps}

\begin{steps}
\item \textbf{Stack the shield and set Mode 1.} The IKS4A1 goes on the Nucleo Arduino
headers. Set the UM3239 Mode 1 positions, which put all sensors on the microcontroller I2C
bus:
\begin{plaincode}
J4: 1-2 and 11-12
J5: 1-2 and 11-12
\end{plaincode}
The IKS5A1 stays in its own bag. The two shields look similar from across the bench, and
mixing them is how the wrong chapter gets built.

\item \textbf{Flash the firmware.} Use an STM32duino X-NUCLEO-IKS4A1 sample, or a sketch
that prints one JSON object per 20 ms:
\begin{plaincode}
{"t":24.1,"rh":41.2,"p":1013.2,"ax":0.02,"ay":0.01,"az":1.00}
\end{plaincode}
One object per line, one line per sample period. Keep the key names stable from the first
capture, because a schema that drifts halfway through a file is worse than a short file.

\item \textbf{Record on the Pi.} Connect the Nucleo USB cable to a Pi USB-A port and write
the stream to a file while you watch it.
\begin{shellcode}
cat /dev/ttyACM0 | tee p07.jsonl
\end{shellcode}

\item \textbf{Prove the boundary.} Run the Pi scanner once, so the empty grid is something
you saw rather than something you were told.
\begin{shellcode}
sudo i2cdetect -y 1
\end{shellcode}
It shows nothing, and it is meant to. These chips live behind the STM32. The addresses
0x6A, 0x1E, 0x44, 0x5C and the 0x38-class temperature part are visible only to the scanner
inside the firmware.

\item \textbf{Compute a running tilt from acc and gyro only.} Magnetometer heading waits on
a hard-iron calibration you actually perform. Do not publish a raw heading and call it
fusion. A tilt that comes from two sensors you trust is worth more than a heading that
comes from three you have not calibrated.

\item \textbf{Test the climate parts with the two stimuli.} Breathe toward the SHT40 and
watch relative humidity move. Hold a palm over the board and watch the STTS22H move by
about 0.5 degrees C. Do them separately: a breath carries both moisture and heat, so doing
both at once tells you less than either alone.

\item \textbf{Leave Qvar for a second session.} The Qvar swipe stays off until the IMU
stream is clean. Then enable the electrode jumpers, in the UM3239 Qvar mode, as its own
session with its own capture file, not as a kitchen sink on top of a stream you are still
validating.
\end{steps}

\subsection*{Acceptance test}

\begin{itemize}
\item The firmware scanner sees 0x6A, 0x1E, 0x44, 0x5C and the 0x38-class addresses.
\item \texttt{sudo i2cdetect -y 1} on the Pi shows an empty grid, and that is a pass.
\item Breathing on the SHT40 moves relative humidity.
\item A palm over the board moves the STTS22H by about 0.5 degrees C.
\item \texttt{p07.jsonl} holds one JSON object per line with the same keys from the first
line to the last.
\end{itemize}

\subsection*{Practices}

Mode 1 is the Linux-friendly setup. Sensor-hub and ISPU modes hide devices behind the
LSM6, which is useful after the JSON is honest and confusing before it. Keep this lab's
numbers in their own log file: P08 measures a different instrument and its schema is not
this one.

Record first and derive later. The file is the artifact, the tilt is a view of it, and a
heading you have not calibrated is neither. When you do run the hard-iron calibration,
record it as its own capture with its own note in the lab book, so the correction can be
recomputed rather than remembered.

\subsection*{Pitfalls}

\begin{itemize}
\item \textbf{Running \texttt{i2cdetect} on the Pi and reading the empty grid as a fault.}
The chips are behind the STM32. The symptom is an hour spent checking jumpers that were
correct, followed by a Mode change that actually does hide devices.
\item \textbf{Choosing a sensor-hub or ISPU mode by accident.} Those modes route parts
through the LSM6, so the firmware scan finds fewer addresses than the shield carries. The
symptom is a JSON object that silently loses keys.
\item \textbf{Publishing a raw magnetometer heading as fusion.} Without a hard-iron
calibration the heading follows the nearest steel object. The symptom is a compass that is
stable and wrong.
\item \textbf{Mixing this schema with P08.} The industrial shield produces different
quantities from different silicon. Two schemas in one file is how a wearable-grade number
ends up labelled industrial.
\item \textbf{Stacking both shields.} The Arduino stack takes one shield. The symptom is an
address clash at 0x6A and a scan that reports the wrong part.
\item \textbf{Enabling Qvar before the IMU stream is clean.} Two new variables at once
means neither is isolated when something looks wrong.
\item \textbf{Breathing on the board to test temperature.} A breath moves humidity and
temperature together. Use the palm for STTS22H and the breath for SHT40, one at a time.
\item \textbf{A schema that drifts.} Adding a key halfway through a capture leaves a file
that no parser reads end to end. Restart the capture instead.
\end{itemize}

\subsection*{Sources}

\begin{itemize}
\item ST X-NUCLEO-IKS4A1 expansion board page,
\url{https://www.st.com/en/ecosystems/x-nucleo-iks4a1.html}
\item ST UM3239, the X-NUCLEO-IKS4A1 user manual, for the J4 and J5 mode tables,
\url{https://www.st.com/resource/en/user_manual/um3239.pdf}
\item STM32duino X-NUCLEO-IKS4A1 library and samples,
\url{https://github.com/stm32duino/X-NUCLEO-IKS4A1}
\item ST LSM6DSV16X and LSM6DSO16IS data sheets,
\url{https://www.st.com/en/mems-and-sensors/inemo-inertial-modules.html}
\item ST LPS22DF pressure sensor and STTS22H temperature sensor data sheets,
\url{https://www.st.com/en/mems-and-sensors/pressure-sensors.html}
\item Sensirion SHT40 humidity and temperature sensor data sheet,
\url{https://sensirion.com/products/catalog/SHT40}
\item ST NUCLEO-H7A3ZI-Q board page,
\url{https://www.st.com/en/evaluation-tools/nucleo-h7a3zi-q.html}
\end{itemize}
